\section{Limitations}
Whilst ADS-B trajectories enable aircraft exhaust plume intersections to be identified, there are limitations regarding usage for exhaust plume modelling. Firstly, identification of exhaust plume interactions is limited to global regions with ADS-B OpenSky Network ground receiver coverage, and therefore global estimates of aircraft exhaust plume intersections are not currently feasible. This limitation includes areas where sustained exhaust plume interactions are likely, such as oceanic crossings, particularly along the North Atlantic Flight Corridor, where densities are typically high \cite{Schlager1997}.

Another significant limitation of relying on solely ADS-B derived-trajectories is the lack of meteorological information. The accurate modelling of exhaust plumes and contrails will of course be reliant on the pre-existing atmospheric conditions, along with an estimate plume drift due to wind conditions \cite{Gruber2018}. Consequently, the study presented in this paper has assumed that the plume follows the aircraft ground track. Future work could employ the meteorological information supplied my Mode-S data, or alternative approaches for estimating wind conditions via Mode-S data \cite{Sun2018}, for the aircraft identified with plume intersections within the ADS-B dataset to generated improved plume modelling. When considering 3D plume modelling, further investigations into variations of reported GPS and barometric altitudes of operational aircraft is required, especially for the discrete `steps' in aircraft altitude observed within ADS-B trajectories.

A final limitation with respect to considering solely ADS-B trajectories for identifying aircraft exhaust plume intersections is a result of the plume size and dispersion being related to aircraft operational parameters, including aircraft size (i.e.\ mass) and flying speeds. Both of these parameters will impact the engine exhaust flow and the wake generated by the aircraft, which are expected to impact the size and dispersion of the aircraft exhaust plume. Future work would aim to address this by coupling ADS-B trajectories with aircraft characteristics, prior to conducting plume modelling runs.